\documentclass[11pt]{article}
\usepackage{explorer}
\newcommand{\ExplorerLanguage}{FR}
\newcommand{\ExplorerNumber}{02}
\newcommand{\ExplorerKind}{Solutions détaillées}
\begin{document}
\ExplorerTitle{Seconde · olympiades}{Le principe des tiroirs}
\ExplorerSection{Exercice 1}
Montrer que parmi 25 personnes, trois au moins sont nées le même mois.\par
Avec au plus deux personnes par mois, il y aurait au plus 24 personnes. Donc un mois en contient au moins trois.\par
\ExplorerSection{Exercice 2}
On choisit six nombres distincts dans \{1,…,10\}. Prouver que deux ont pour somme 11.\par
Les cinq tiroirs sont \{1,10\}, \{2,9\}, \{3,8\}, \{4,7\}, \{5,6\}. Six choix occupent entièrement au moins une paire.\par
\ExplorerSection{Exercice 3}
Prouver que parmi huit entiers, deux ont une différence divisible par 7.\par
Il y a sept restes possibles modulo 7. Deux des huit entiers ont le même reste, donc leur différence est multiple de 7.\par
\ExplorerSection{Exercice 4}
Cinq points sont dans un carré de côté 2, bord compris. Montrer que deux sont à distance au plus \(\sqrt{2}\).\par
Partager le carré en quatre carrés de côté 1, en affectant chaque point des frontières à un seul carré. Deux points appartiennent au même petit carré ; sa diagonale vaut \(\sqrt{2}\).\par
\ExplorerSection{Problème : raisonnement complet}
On choisit 11 entiers distincts entre 1 et 20. Écrire chaque entier sous la forme 2ᵏm avec m impair et \(k \geq  0\). Combien de valeurs de m sont possibles ? Prouver que l’un des entiers choisis divise un autre. Montrer qu’avec seulement dix entiers, la conclusion peut être fausse.\par
La décomposition s’obtient en retirant les facteurs 2. Les dix parties impaires possibles sont 1,3,…,19. Deux nombres choisis partagent donc la même partie impaire : 2ᵃm et 2ᵇm, avec \(a < b\) puisqu’ils sont distincts. Le second vaut 2ᵇ⁻ᵃ fois le premier. Pour dix nombres, choisir 11,12,…,20 : le double de chacun dépasse 20, donc aucun n’en divise un autre distinct.\par
\end{document}
